% =================================================
% A 册补充专题：性侵犯与强奸（认知与应对）
% 现代循证撰写，紧跟「沟通与同意」章
% =================================================

\chapter{性侵犯与强奸：认知与应对}

\section{先正名：这不是"你的错"，也不是"她没拒绝够大声"}
性侵犯（sexual assault / rape）的核心是\textbf{同意的缺失}，而不是"用了什么方式""在哪里发生"。
只要一方没有\textbf{自愿、清醒、可撤回}的同意，就是侵犯——无论发生在婚姻内、恋爱中、职场、校园，还是陌生人之间。
把它单独成章，是因为它\textbf{不是女性专利，也不是"小事"，更不是某个性行为的变体}，而是一种以权力与暴力为底色的侵犯。
理解它的第一课，就是把"受害者有罪论"请出去。

\section{常见误解（请逐条拆开）}
\begin{tabularx}{\textwidth}{l X}
\toprule
\textbf{误解} & \textbf{事实} \\
\midrule
"他没穿裙子 / 没喝醉，就不会被侵犯" & 是否穿着、是否饮酒，都与"谁在施加、谁没同意"无关。 \\
"都结婚了，谈不上强奸" & 婚内强奸在多数国家（含中国刑法修订后）已被承认为犯罪；同意不是结婚证附带的。 \\
"她没大声喊'不'，就是默许" & 沉默、僵硬（freezing）、事后才说出"不"，都可能是同意缺失的表现；同意须是\textbf{主动且持续}的。 \\
"隔了很久才说，是编的" & 创伤后记忆常是非线性的；羞耻、恐惧、关系压力都会让人延迟求助，延迟不等于不真。 \\
"只有暴力才算" & 胁迫、利用职权/信息差、事后无法撤回的同意，都算侵犯，未必有淤青。 \\
\bottomrule
\end{tabularx}

\section{当下该怎么做（第一反应清单）}
目标是\textbf{先保证安全，再处理后续}，顺序可随情况灵活调整：
\begin{enumerate}
  \item \textbf{到安全的地方}：离开现场，去能信任的人身边，或拨打 110 / 当地求助热线。
  \item \textbf{不要急着"洗掉证据"}：在未就医取证前，尽量避免洗澡、刷牙、换衣、排便——这些会冲掉唾液、精液、皮屑、纤维等物证。若不得不清洗，尽量用干净容器保存衣物。
  \item \textbf{尽快就医（最好 72 小时内）}：到\textbf{急诊或妇科/男科}，说明"可能需要性侵取证（rape kit）"。医生会检查伤情、取物证、并评估 STI 与妊娠风险。
  \item \textbf{取证与报案可分开做}：可以先体检、保留证据，之后再决定是否走法律程序；证据保存期通常有限，但\textbf{身体安全永远排在第一}。
  \item \textbf{记录与见证}：条件允许时，记下时间、地点、人物特征、对话，保存聊天记录/监控线索；有人陪同更好。
\end{enumerate}

\section{健康后果：可查、可治、可防}
性侵不只是"那一次"，它带来一连串\textbf{可处理}的医学问题，别自己扛：
\begin{itemize}
  \item \textbf{性传播感染（STI）暴露}：HIV、梅毒、淋病、HPV 等风险上升。就医后可评估是否需要\textbf{HIV 暴露后预防（PEP，72 小时内开始最有效}）、抗生素预防及后续随访检测（通常 2--4 周、3 个月、6 个月复测）。
  \item \textbf{意外妊娠}：无保护性交后，紧急避孕（口服紧急避孕药 72 小时内，或铜宫内节育器 5 天内）成功率较高；后续可评估常规避孕方案。
  \item \textbf{生殖器 / 泌尿道损伤}：撕裂、血尿、疼痛，需检查处理；女性创伤细节见\textbf{本系列 D 册“女性性健康”}相关章节。
  \item \textbf{慢性盆腔痛 / 性交痛（dyspareunia）}：若持续，常与创伤后肌肉紧张、心理因素相关，可治疗（见 B 册"常见困惑深挖"、D 册相关专题）。
\end{itemize}

\section{长期心理影响：可以恢复，不必"想开点"}
\begin{itemize}
  \item 常见反应：闪回、失眠、回避、情绪麻木、信任感下降，重者出现\textbf{创伤后应激（PTSD}）。这些是\textbf{正常人对异常事件}的应激，不是"脆弱"。
  \item 求助通道：心理科 / 精神科、心理咨询（推荐\textbf{创伤聚焦认知行为 TF-CBT} 或\textbf{EMDR}），以及本地妇女维权 / 心理援助热线。
  \item 关系修复（如与伴侣）：以\textbf{自愿、可撤回、循序渐进}为前提，不强迫"复原到从前"；必要时接受伴侣共同咨询。
  \item 一句给当事人的话：\textbf{你不需要为别人的暴力负责；你有权按自己的节奏处理，而不是被别人催着"放下"。}
\end{itemize}

\section{男性也是受害者（常被沉默的那一半）}
\begin{itemize}
  \item 统计上女性受害率更高，但\textbf{男性、非顺性别者同样会发生}；男性受害者常被"男人该扛住"的刻板印象耽误求助。
  \item 男性同样面临：STI / 妊娠（对伴侣）风险、生殖道损伤、PTSD，以及更强的羞耻与"不被相信"感。
  \item 处理路径与上文一致：安全→就医取证→STI/妊娠评估→心理支持；男性可到\textbf{男科或泌尿外科}说明情况，不必因"面子"硬扛。详见\textbf{本系列 C 册“男性性健康”}。
\end{itemize}

\section{预防与保护：把"同意"做成习惯}
\begin{itemize}
  \item \textbf{事前}：了解并尊重彼此的边界；饮酒后决策力下降，重要决定尽量在清醒时沟通好。
  \item \textbf{事中}：把"可以吗？还想继续吗？"当成自然的询问，而不是"质问"；\textbf{同意可以撤回}，任何时刻说"停"都有效。
  \item \textbf{环境}：公共/半公共场合的性行为注意隐私与照明；夜间出行尽量结伴、保持联络；职场/师生等权力不对等关系中，对"难以拒绝"的邀请保持警觉。
  \item \textbf{对他人}：被求助时，第一反应是"信 + 陪伴 + 不追问细节"，而不是"你怎么不早点说 / 你当时怎么没跑"。
\end{itemize}

\section{一页要点}
性侵犯 = 同意缺失，\textbf{不是受害者的错}。当下先保安全、再谈取证与就医（72 小时内就医 + PEP 窗口）；
健康后果（STI/妊娠/损伤）可查可治，别硬扛；心理影响可以恢复，按自己的节奏来。
男性同样是受害者，路径相同。预防的核心是\textbf{把"同意"做成日常习惯}，而不是事后追查。
相关女性专科细节见 D 册，男性专科细节见 C 册。


\chapter{附录：关键术语对照与权威来源}

\section{关键术语首现对照（便于跨册速查）}
以下缩写/术语在五册中反复出现，首次出现时请按下表理解，之后不再重复解释：

\begin{tabularx}{\textwidth}{l X X}
\toprule
\textbf{缩写/术语} & \textbf{中文全称} & \textbf{一句话说明} \\
\midrule
ED & 勃起功能障碍（erectile dysfunction） & 持续无法获得或维持足以性交的勃起，$\ge$3 个月；常与心血管、内分泌、心理相关。 \\
PEP & HIV 暴露后预防（post-exposure prophylaxis） & 疑似 HIV 暴露后 72 小时内开始、连服 28 天抗逆转录病毒药，可大幅降低感染率。 \\
PrEP & HIV 暴露前预防 & 高危人群每日服抗逆转录病毒药，事前降低感染风险。 \\
STI & 性传播感染（sexually transmitted infection） & 泛指 HIV/梅毒/淋病/HPV/衣原体等，感染（infection）比"病（disease）"更准确。 \\
HPV & 人乳头瘤病毒 & 分高危型（致癌：16/18 等）与低危型（疣：6/11）；疫苗可防主流高危型。 \\
SSRI & 选择性 5-羟色胺再摄取抑制剂 & 一线抗抑郁药；常见副作用为性欲下降、射精延迟、女性高潮困难。 \\
PDE5i & 5 型磷酸二酯酶抑制剂（万艾可/希爱力等） & 一线 ED 药物；需性刺激才起效；禁与硝酸酯类同服。 \\
TF-CBT & 创伤聚焦认知行为治疗 & PTSD/性侵后抑郁的一线心理治疗，聚焦"创伤记忆 + 认知 + 回避"。 \\
EMDR & 眼动脱敏与再加工 & 通过双侧刺激处理创伤记忆，对 PTSD 有较强证据。 \\
Kegel & 凯格尔运动 & 盆底肌收缩训练；男女皆适用，需正确"上提"方向（见 B 册“呼吸与盆底肌（Kegel）技巧”）。 \\
edging & 边缘控制 / start--stop & 接近高潮时退后、停留再升，延时长 + 增强度（见 B 册“高潮引导与边缘控制”）。 \\
G-spot & "G 点" & 阴道前壁 2--3cm 处的海绵体组织，个体差异大，并非人人有明显反应。 \\
SSC & 安全·理性·知情同意（Safe, Sane, Consensual） & BDSM 圈通行原则，所有玩法的前提。 \\
S\&M & 施虐（sadism）·受虐（masochism） & BDSM 中的两大核心动力，可单独或结合。 \\
D/s & Dominant / submissive & 支配者 / 服从者角色，与性别无关。 \\
aftercare & 事后关怀 & 高强度玩法（含 orgasm）后的情绪/身体照护，BDSM 与一般性都适用。 \\
\bottomrule
\end{tabularx}

\section{权威来源清单（正文引用与延伸阅读）}
本系列正文的事实与数据主要来自以下机构/共识/教材，按主题分类：

\begin{itemize}
  \item \textbf{全球性教育标准（ICSE）}：\href{https://www.wego.org}{wego.org}（WHO/UNESCO 联合维护），用于本系列"性是什么、性反应周期、同意与沟通"章节。
  \item \textbf{WHO / 联合国人口基金（UNFPA）}：生殖健康与性健康定义、全球 STI 指南、HPV 疫苗接种策略。
  \item \textbf{CDC（美国疾控中心）}：STI 诊疗指南（HIV/梅毒/淋病/衣原体/HPV）、PEP/PrEP 推荐、紧急避孕窗口。
  \item \textbf{EAU（欧洲泌尿外科学会）}：ED 诊疗指南、早泄（PE）共识、前列腺健康。
  \item \textbf{AUA / NICE（英国国家卫生与临床优化研究所）}：ED 与 PE 的一线治疗路径、盆底肌康复。
  \item \textbf{ACOG（美国妇产科医师学会）}：经期/孕期/更年期性生活、避孕方法可靠性。
  \item \textbf{DSM-5（精神障碍诊断与统计手册 第5版）}：性功能障碍分类（欲/兴/高潮/痛等）。
  \item \textbf{NHS（英国国家医疗服务体系）}：面向大众的性健康、PEP/PrEP、避孕选择。
  \item \textbf{中国}：《中国人类精子库技术规范》、国家卫健委 STI 诊疗指南、《中国女性健康指南》相关章节。
  \item \textbf{经典教材}：Human Sexual Response Cycle（Masters \& Johnson, 1966）、Koopman \& Berman（性反应模型更新）、Kaplan（三阶段模型）。
\end{itemize}

\textbf{使用提示}：本系列面向大众，非处方或诊断依据；具体诊疗请结合\textbf{主诊医生}建议。正文中出现"建议就医/复测"处，均为循证节点，不是营销话术。


